\section*{Exercises about Bases}

\exercise{6}
\begin{xrcs}
  Claim: If $v_1, v_2, v_3, v_4$ is a basis of $V$, then
  \begin{equation}
    v_1 + v_2, \; v_2 + v_3, \; v_3 + v_4, \; v_4 \text{ is also a basis of } V.
  \end{equation}

  \begin{xprf}
    \StepOne We first show linear independence. Let $a_1, a_2, a_3, a_4 \in \myF$ such that
    \begin{equation}
      a_1 (v_1+v_2) + a_2(v_2+v_3) + a_3(v_3+v_4) + a_4 v_4 = 0.
    \end{equation}
    Expanding this expression gives:
    \[
      a_1 v_1 + (a_1 + a_2) v_2 + (a_2 + a_3) v_3 + (a_3 + a_4) v_4 = 0.
    \]
    Since $v_1, \ldots, v_4$ is a basis (hence linearly independent), the coefficients must vanish:
    \[
      a_1 = 0, \quad a_1 + a_2 = 0, \quad a_2 + a_3 = 0, \quad a_3 + a_4 = 0.
    \]
    Since $a_1 = 0$, back-substitution immediately yields $a_2 = 0, a_3 = 0, a_4 = 0$.
    Thus, the list $v_1 + v_2, v_2 + v_3, v_3+v_4, v_4$ is linearly independent.

    \StepTwo We now show that the list spans $V$. Let $u \in V$. Since $v_1, \dots, v_4$ is a basis, there exist scalars $b_1, \dots, b_4 \in \myF$ such that
    \begin{equation}
      \label{eq: equation of u}
      u = b_1 v_1 + b_2 v_2 + b_3 v_3 + b_4 v_4.
    \end{equation}
    We can express $u$ as a linear combination of the new list by solving for coefficients $a_1, \dots, a_4$:
    \[
      a_1 = b_1, \quad a_2 = b_2 - b_1, \quad a_3 = b_3 - b_2 + b_1, \quad a_4 = b_4 - b_3 + b_2 - b_1.
    \]
    Substituting these gives
    \[
      u = b_1 (v_1+v_2) + (b_2-b_1)(v_2+v_3) + (b_3 - b_2 + b_1)(v_3 + v_4) + (b_4 - b_3 + b_2 - b_1) v_4.
    \]
    Thus, the list spans $V$. Since it is also linearly independent, it is a basis of $V$.
  \end{xprf}
\end{xrcs}

\exercise{8}
\begin{xrcs}
  Prove or give a counterexample: If $v_1, v_2, v_3, v_4$ is a basis of $V$ and $U \subseteq V$ is a subspace such that $v_1, v_2 \in U$ and $v_3 \notin U$ and $v_4 \notin U$, then $v_1, v_2$ is a basis of $U$.

  \begin{xsol}
    The statement is false.

    \prooffont{Counterexample:} In $V = \real^4$, let $v_1, v_2, v_3, v_4$ be the standard basis. Define
    \begin{equation}
      U := \myspan{v_1, v_2, v_3 + v_4}.
    \end{equation}
    Then $v_1, v_2 \in U$, while $v_3 \notin U$ and $v_4 \notin U$.
    However, the three vectors $v_1, v_2, v_3+v_4$ are linearly independent, so $\dim U = 3$.
    Therefore, the list $v_1, v_2$ (which has length $2$) cannot span $U$ and is not a basis of $U$.
  \end{xsol}

  \begin{failedattempt}
    Any vector in $U$ can be represented as a unique linear combination of $v_1, v_2, v_3, v_4$ since $U \subseteq V$. Let $u\in U$ such that $u=a_1 v_1 + a_2 v_2 + a_3 v_3 + a_4 v_4$. In this case, $a_3$ must equal $0$, because otherwise, due to closedness under addition and scalar multiplication, we could subtract $a_1 v_1 + a_2 v_2+a_4v_4$ from $u$ and multiply by $\frac{1}{a_3}$ to get $v_3$, which is not in $U$. The same goes for $a_4$, which must be $0$ as well. So every vector in $U$ is of the form $u=a_1 v_1 + a_2 v_2$, since $v_1, v_2 \in U$ and $v_3, v_4 \notin U$. Since every basis has the same length and $v_1$ and $v_2$ are linearly independent, the list $v_1, v_2$ forms a basis of $U$.
  \end{failedattempt}

  \begin{ainote}[Where it breaks]
    The sentence \qt{we could subtract $a_1 v_1 + a_2 v_2 + a_4 v_4$ from $u$} is the faulty step. Subtracting is only allowed inside $U$ if the vector being subtracted is itself in $U$, and $a_4 v_4$ need not be: $v_4 \notin U$ is given. So the argument quietly assumes that each component $a_k v_k$ of $u$ lies in $U$ separately, which is precisely what a subspace does \emph{not} guarantee. A subspace is closed under sums, but a sum landing in $U$ says nothing about its individual summands.

    The counterexample above is exactly this loophole made concrete. For $u := v_3 + v_4 \in U$ the coefficients are $a_3 = a_4 = 1 \neq 0$, so the attempt's claim \qt{$a_3$ must equal $0$} is already false, even though $v_3 \notin U$ and $v_4 \notin U$ both hold as required.

    A repair is possible, but only under a stronger hypothesis: if one assumes $\myspan{v_3,v_4} \cap U = \{0\}$, then the components really do split off and $v_1,v_2$ is a basis of $U$. The exercise deliberately does not give you that.
  \end{ainote}
\end{xrcs}

\exercise{9}
\begin{xrcs}
  Suppose $v_1, \dots, v_m$ is a list of vectors in $V$. For $k \in \{1, \ldots, m\},$ let
  \begin{equation}
    w_k := v_1 + \cdots + v_k.
  \end{equation}
  Show that $v_1, \ldots, v_m$ is a basis of $V \iff w_1, \ldots, w_m$ is a basis of $V$.

  \begin{xprf}
    By Exercise 2A.4, $\myspan{w_1, \dots, w_m} = \myspan{v_1, \dots, v_m}$.
    Thus, $w_1, \dots, w_m$ spans $V$ if and only if $v_1, \dots, v_m$ spans $V$.

    Now we check linear independence:
    \begin{itemize}
      \item Suppose $v_1, \dots, v_m$ is linearly independent. If $c_1 w_1 + \cdots + c_m w_m = 0$, expanding in terms of the $v$'s yields:
      \[
        (c_1 + c_2 + \cdots + c_m) v_1 + (c_2 + \cdots + c_m) v_2 + \cdots + c_m v_m = 0.
      \]
      By linear independence of $v_1, \dots, v_m$, each coefficient is zero:
      \[
        c_m = 0 \implies c_{m-1} = 0 \implies \cdots \implies c_1 = 0.
      \]
      Thus, $w_1, \dots, w_m$ is linearly independent.

      \item Conversely, suppose $w_1, \dots, w_m$ is linearly independent. If $a_1 v_1 + \cdots + a_m v_m = 0$, substituting $v_1 = w_1$ and $v_k = w_k - w_{k-1}$ for $k \ge 2$ yields:
      \[
        (a_1 - a_2) w_1 + (a_2 - a_3) w_2 + \cdots + (a_{m-1} - a_m) w_{m-1} + a_m w_m = 0.
      \]
      By linear independence of $w_1, \dots, w_m$:
      \[
        a_m = 0 \implies a_{m-1} = 0 \implies \cdots \implies a_1 = 0.
      \]
      Thus, $v_1, \dots, v_m$ is linearly independent.
    \end{itemize}
    Since both spanning and linear independence are equivalent, $v_1, \dots, v_m$ is a basis of $V \iff w_1, \dots, w_m$ is a basis of $V$.
  \end{xprf}
\end{xrcs}
